%% -------------Packages
% margin
\usepackage[margin=0.9in]{geometry}

% spelling
\usepackage[english]{babel}

% headers and footers
\usepackage[twoside]{fancyhdr}

% sectioning
\usepackage{titlesec}

% table of content
\usepackage{titletoc}

% line spacing
\usepackage{setspace}

% font settings for captions
\usepackage[font=small,textfont=sl]{caption}

% bibliography
\usepackage[square,numbers]{natbib}
\bibliographystyle{plain}

% For spaces after my commands
\usepackage{xspace}

% make references clickable 
\usepackage[]{hyperref}
\def\manuscriptTitle{Reactive Programming Language Applied to Automotive Functional Specification and Model-Based Design\xspace}
\hypersetup{
  pdftitle={\manuscriptTitle},
  % pdfpagemode=FullScreen,
}

% for maths symbols
\usepackage{stmaryrd, amssymb, amsthm}

% for figures
\usepackage{graphicx}
\graphicspath{ {../images/} }

% for the `H' option (hear)
\usepackage{float}

% for BNF grammar syntax 
\usepackage{simplebnf}

% for bold math symbols
\usepackage{bm}

% logical terms
\usepackage{mathpartir}
\makeatletter
\let\originferrule\inferrule
\DeclareDocumentCommand \inferrule { s O {} m m o }{%
  \IfBooleanTF{#1}%
  {%
    \mpr@inferstar[#2]{#3}{#4}%
  }{%
    \mpr@inferrule[#2]{#3}{#4}%
  }%
  \IfValueT{#2}%
  {%
    \my@name@inferrule{#2}%
  }%
}
\NewDocumentCommand \my@name@inferrule { m }{%
  \def\@currentlabelname{\ensuremath{#1}}%
}
\makeatother

% boxes
\usepackage{tcolorbox}
% New boxes
\newtcolorbox[auto counter]{syntaxbox}[2][]{%
  top=0mm,
  bottom=0mm,
  left=0mm,
  colback=blue!10,
  colframe=blue!50!black,
  title=Syntax.~\thetcbcounter: #2,
  #1
}
\newtcolorbox[auto counter]{semanticsbox}[2][]{%
  top=0mm,
  bottom=0mm,
  left=0mm,
  colback=red!10,
  colframe=red!50!black,
  title=Semantics.~\thetcbcounter: #2,
  #1
}
\newtcolorbox[auto counter]{examplebox}[2][]{%
  top=0mm,
  bottom=0mm,
  left=0mm,
  colback=green!10,
  colframe=green!50!black,
  title=#2,
  #1
}

% latin modern fonts
\usepackage{lmodern}

% multiple columns
\usepackage{multicol}

% automatic line breaks in tables + long tables
\usepackage{tabularx, longtable, booktabs}

% combining cells in tabular
\usepackage{multirow}

% listings
\usepackage{listings, listings-extended}
\lstset{
inputpath=../codes,
literate=
  {à}{{\`a}}1
{é}{{\'e}}1
{è}{{\`e}}1
{ê}{{\^e}}1
{î}{{\^i}}1
{ô}{{\^o}}1
{û}{{\^u}}1
{ç}{{\c{c}}}1
}

% For colors
\usepackage[dvipsnames]{xcolor}

% For conditional macros with optional arguments
\usepackage{ifthen}

% For wrapping figures
\usepackage{wrapfig,caption,subcaption}
\setlength\intextsep{5pt}

% For graphics and drawings
\usepackage{tikz}
\usetikzlibrary{arrows.meta,shapes,arrows,calc,positioning,fit}
% Define block styles
\tikzstyle{block} = [
rectangle,
rounded corners,
draw,
text width=2cm,
text centered,
minimum height=1cm
]
\tikzstyle{arrow}=[-{Latex[length=2mm, width=2mm]}]
\tikzstyle{crate_border}=[
dotted,
thick,
rounded corners=5mm,
inner sep=5mm,
minimum height=2.5cm,
minimum width=3.5cm
]
\tikzstyle{process_block}=[block, fill=blue!10]
\tikzstyle{data_block}=[block, fill=green!10]
\tikzstyle{label_block}=[
block,
rounded corners=5mm,
draw = white,
nearly opaque,
fill=red!10,
text width=9cm,
]
\tikzstyle{line} = [draw, -latex']

% % For numbering lines
% \usepackage{lineno}
% \linenumbers

% for reference
\usepackage{cleveref}
